\section{Related Work}
\label{sec:related}

\paragraph{Inference-time scaling for diffusion models.}
Inference-time scaling improves sample quality by spending additional
computation at test time without retraining the
model~\citep{ma2024inference,ramesh2025nts}.
\citet{ma2024inference} systematize this direction by proposing a framework
organized around two design axes---the \emph{verifier} (any external quality
metric) and the \emph{search algorithm} (best-of-$N$, zero-order search, or
path search)---but apply each algorithm with a fixed, per-timestep-uniform
strategy and do not learn which timesteps benefit most from search.
Our work extends this line of research by focusing on the complementary
question of \emph{how to allocate} search compute across the trajectory.

\paragraph{Noise search and trajectory refinement.}
Several methods perform local search over noise variables during diffusion
sampling.
A growing body of work establishes that noise quality varies markedly
across timesteps and samples: \citet{qi2024noises} show that not all initial
noises are equally suitable for generation and propose noise selection via
inversion stability, while \citet{zhou2025golden} learn a noise prompt network
that transforms random Gaussian noise into ``golden noise'' optimized for
text--image alignment.
These results motivate per-step noise refinement at inference time.
Best-of-$N$ methods sample multiple full trajectories and select the best final
output, but do not refine noise within a trajectory~\citep{ma2024inference}.
\citet{tang2024dno} formulate noise refinement as continuous optimization
(Direct Noise Optimization), applying gradient- or zero-order updates at every
timestep, again with uniform per-step effort.
\citet{ramesh2025nts} cast diffusion as an MDP and relax it to a sequence of
independent contextual bandits, proposing an $\epsilon$-greedy algorithm that
implicitly adapts exploration vs.\ exploitation across timesteps; they note
in an appendix that varying $K_t$ across timesteps can reduce NFE by half
with minimal quality loss, but leave the design of an adaptive $K_t$ scheduler
as future work.
\citet{kim2025rbf} introduce Rollover Budget Forcing (RBF), which rolls over
unused per-step budget to subsequent timesteps---a purely online, reactive
mechanism without prior knowledge of timestep importance.
RBF is developed and validated exclusively on flow models (FLUX),
relying on ODE-to-SDE conversion for particle sampling; its applicability
to DDPM or SDE-based diffusion models remains unvalidated.
\citet{zhang2025tscend} propose T-SCEND, a two-phase strategy that
applies best-of-$N$ in early steps and MCTS in late steps; however, the
phase split is fixed and hand-designed rather than data-driven.
More recently, \citet{zhang2025classical} combine local Langevin MCMC with
global tree search (BFS/DFS) for inference-time control, and
\citet{lee2025abcd} introduce Adaptive Bi-directional Cyclic Diffusion (ABCD),
which dynamically adjusts exploration depth per instance.
More recently, \citet{ttsnap2025} propose TTSnap, which prunes low-quality
noise candidates early without full denoising, reallocating saved compute to
explore more candidates; however, this operates at the \emph{candidate} level
within a single timestep rather than scheduling budget \emph{across} timesteps.
In the particle-based setting, \citet{singhal2025fk} introduce a Feynman-Kac
steering framework that allocates non-uniform particle budgets across denoising
stages, and \citet{kim2025psi} propose $\Psi$-Sampler for SMC-based
inference-time alignment with improved initial particle sampling.
\citet{eyring2025hypernetworks} take an orthogonal approach by amortizing
noise search into a learned Noise Hypernetwork, replacing test-time
optimization with a single forward pass.
While these concurrent methods introduce various forms of non-uniform
allocation, they are either purely reactive (RBF, ABCD), use fixed phase
splits (T-SCEND), operate at the candidate rather than timestep level (TTSnap),
or use SMC resampling rather than search-based allocation (FK Steering,
$\Psi$-Sampler).
Our framework explicitly combines \emph{offline profiling}---learning which
timesteps are inherently more amenable to search---with \emph{online
control} that adapts per-step effort to each prompt, achieving non-uniform
budget allocation with strict NFE guarantees.

\paragraph{Adaptive compute allocation.}
Adaptive computation has been studied extensively in other domains, including
early exiting in neural networks, adaptive depth inference, and budgeted
optimization.
In diffusion models, \citet{zhang2025adadiff} propose AdaDiff for adaptive
denoising step selection conditioned on input complexity,
\citet{ye2025tpdm} learn a per-instance time prediction module via RL that
decides which timesteps to visit (adapting the \emph{schedule} rather than the
per-step \emph{search budget}), and
\citet{yeh2025demons} introduce Sampling Demons, a training-free alignment
method that reweights per-step noise importance to steer generation toward
high-reward regions.
However, these techniques primarily modify the base sampler or noise
distribution rather than allocating additional \emph{search} computation
across timesteps.
\citet{sundaresha2025vt} propose the Verifier-Threshold method for more
efficient non-uniform NFE allocation, demonstrating 2--4$\times$ speedups
over RBF; however, their approach is also purely online and does not use
offline profiling.
Our approach is complementary to sampler design: we keep the underlying sampler
fixed and focus on how to schedule inference-time search around it, combining
a learned offline allocation with online gain-and-variance feedback.
